\chapter{Docker: Images, Containers, Compose}
\label{ch:docker}

\section{What a container actually is}

A container is not a small virtual machine. It is an ordinary Linux
process, isolated by two kernel features: \emph{namespaces} (its own view
of processes, network, filesystem mounts, hostname) and \emph{cgroups}
(limits on CPU and memory). No guest OS, no hypervisor --- which is why a
container starts in milliseconds and why the JVM inside one must be told
about its memory limit (Chapter~17's \code{MaxRAMPercentage}) --- from
inside, it can see the \emph{host's} RAM.

An \emph{image} is the packaged filesystem plus metadata (entrypoint,
env, exposed ports), built as a stack of immutable \emph{layers}. Layers
are content-addressed and shared: fifty images \code{FROM
eclipse-temurin:25-jre-alpine} store the JRE layer once. A
\emph{registry} (Docker Hub, GHCR, the project's \code{registry:2})
stores images; \code{docker push}/\code{pull} move only layers the other
side lacks.

\section{The project's Dockerfile, annotated}

Every service carries the same five lines:

\begin{lstlisting}[language=dockerfile]
FROM eclipse-temurin:25-jre-alpine    (*@\co{1}@*)
WORKDIR /app
COPY target/*.jar app.jar             (*@\co{2}@*)
EXPOSE 8082                           (*@\co{3}@*)
ENTRYPOINT ["java", "-jar", "app.jar"](*@\co{4}@*)
\end{lstlisting}

\begin{description}
  \item[\co{1}] Base image: a JRE (not JDK --- no compiler needed at
  runtime) on Alpine Linux. Smaller base = smaller attack surface and
  faster pulls.
  \item[\co{2}] Copies the Maven-built fat jar. Subtlety: this
  Dockerfile requires \code{mvn package} to have run \emph{first} ---
  the build happens outside. The consequence: one \code{COPY} of a
  \(\sim\)60\,MB jar makes one big layer, so \emph{every} build re-pushes
  everything even for a one-line change. Jib (Chapter~13) exists to fix
  exactly this.
  \item[\co{3}] Documentation, not action --- it opens nothing.
  Publishing ports happens at run time.
  \item[\co{4}] Exec-form entrypoint: java is PID~1, receives signals
  directly, so \code{docker stop} triggers a clean Spring shutdown.
\end{description}

\section{Networking and Compose}

Containers on a Docker \emph{bridge network} reach each other by
\emph{service name} --- Docker runs a DNS server answering
\code{postgres-course}, \code{kafka}, \code{minio}. \code{localhost},
inside a container, is \emph{that container} --- the number-one beginner
trap, and the reason every dev-default URL in the configs needed a
per-environment override. Ports are only reachable from the host when
\emph{published} (\code{-p 5434:5432}).

Compose turns a fleet of \code{docker run} commands into one declarative
file --- the project's \code{docker-compose.yml} is the whole dev stack:
four datastores, Kafka, Zipkin, Maildev, and the seven services, with
three idioms worth naming:

\begin{lstlisting}[language=yaml]
  discovery-server:
    depends_on:
      config-server:
        condition: service_healthy    (*@\co{1}@*)
  config-server:
    healthcheck:
      test: ["CMD", "curl", "-f", "http://localhost:8888/actuator/health"] (*@\co{2}@*)
  postgres-course:
    ports:
      - "5434:5432"                   (*@\co{3}@*)
    volumes:
      - postgres-course-data:/var/lib/postgresql/data
\end{lstlisting}

\co{1}+\co{2} encode startup order --- and not merely ``started'' but
``answering its health endpoint,'' which is what makes the ordering
real. \co{3} maps distinct host ports (5433, 5434) onto the two
Postgres containers' identical internal 5432 --- names are per-network,
host ports are global. The named volume keeps data across container
recreation; Chapter~19 is the Kubernetes version of the same need.

\begin{decision}[why Compose survives alongside Kubernetes]
The cluster is the deployment target, but Compose remains the inner
loop: one command, laptop-local, debug ports mapped (every service opens
a JDWP port in Compose --- attach the IDE debugger to 5082 and step
through course-service). Kubernetes buys self-healing and rolling
deploys at the cost of ceremony; a developer iterating on business logic
wants neither cost. Two environments, one image contract between them.
\end{decision}

\begin{pitfall}[works in Compose, dies in Kubernetes]
Symptom: the stack is green locally; the same image crash-loops in the
cluster. The usual suspects, all environmental: a URL still pointing at
a Compose service name that has a different Service name in k8s; a
\code{depends\_on} ordering the cluster does not honor (Kubernetes has
no startup ordering --- Chapter~15); or memory limits present in the
cluster but absent locally, turning a fat JVM into an OOMKill. The
image is identical --- diff the \emph{environment}, not the code:
\code{docker inspect} vs.\ \code{kubectl describe pod}.
\end{pitfall}

\begin{quiz}
\begin{enumerate}
  \item A container and a VM both isolate. Name two mechanisms the
  container uses, and two consequences of not having a guest kernel.
  \item Why does \code{localhost:5432} fail inside a container that
  works on your laptop, and what two different fixes do Compose and
  Kubernetes offer?
  \item Both Postgres containers listen on 5432 internally without
  conflict --- explain. Where do conflicts become possible again?
  \item What makes \code{condition: service\_healthy} stronger than
  plain \code{depends\_on}, and which project file supplies the health
  signal it waits for?
\end{enumerate}
\end{quiz}
